% Ad Familiares 11.27 --- headnote
% ad-familiares-11-27, 30 August 44 BC, Tusculum

\textit{Cicero to C. Matius, from the Tusculan villa,
between 23 and 30 August 44 BC --- Perseus dateline
\textit{Scr. in Tusculano inter x et iii K. Sept. a. 710
(44)}. C. Matius was an old equestrian friend of Cicero's
and a devoted personal friend of Julius Caesar. After the
Ides, Matius had taken on the management of the
\textit{ludi Victoriae Caesaris} that Octavian put on in
late July, and was widely reported to be the leading voice
of the Caesarian remnant in private conversation in the
city. The complaint to which the letter answers had been
brought from Matius by Trebatius: that Cicero --- or at
any rate Cicero's circle --- was treating Matius as having
betrayed the friendship by aligning with the heirs of the
murdered dictator, and was reporting that Matius had voted
for one of Antony's bills.}

\textit{The letter is one of the major documents in the
corpus on the ethics of friendship under tyranny. Cicero
defends his amicitia with Matius across the whole arc of
their thirty years' acquaintance --- their early
familiarity, Matius's offices on his behalf in 49, the
hurried visits and consolations of the civil-war
years --- and then turns to the present complaint with a
piece of distinctively Ciceronian casuistry: there are two
ways the defence of Matius's conduct can be argued, and
he, Cicero, is accustomed to take the more generous of
them. He insists on the falsity of the rumour about the
vote; he praises Matius's loyalty to a dead friend; but
he also makes very plain, in the eighth section, his own
view that ``the liberty of one's country is to be set
above the life of a friend,'' and that Caesar ``was a
king.''}

\textit{The letter pairs with Matius's celebrated reply
(11.28). Together the two are the closest surviving
exchange we have between a constitutionalist and an
unrepentant Caesarian in the autumn of 44, and the closest
Cicero anywhere comes to writing out, in a letter, the
philosophical case he was simultaneously making in
\textit{De Amicitia} and \textit{De Officiis} for an
\textit{amicitia} bounded by \textit{honestum}. The
Latinity is high-rhetorical throughout --- anaphora at the
heads of the recollections of 49 and 48, the controlled
contrast of \textit{vetustas} and \textit{amor} at the
opening of $\S$~2, the doubled defence-structure of
$\S$~7 (``some things I deny, others I defend''), and the
sustained \textit{utramque partem} disputation at the
close.}
